\documentclass[10pt,fleqn]{article}
\usepackage[utf8]{inputenc}
\usepackage[T1]{fontenc}
\usepackage{fullpage}
\usepackage{amsmath,amssymb,amsthm,mathtools}
\usepackage{mathrsfs} % mathscr
\usepackage{color}
\usepackage[paperwidth=160mm,paperheight=90mm,margin=10mm]{geometry}

% multicol
\usepackage{multicol}
\usepackage{array}

% graphics
\usepackage{graphicx}
\DeclareGraphicsRule{*}{mps}{*}{}

\usepackage{hyperref}
\hypersetup{
  colorlinks=true
}

\usepackage{mathpazo}
\usepackage{helvet}
%\usepackage{courier}
\usepackage{inconsolata}

\usepackage{titling}
\newcommand{\subtitle}[1]{%
  \posttitle{%
    \par\end{center}
    \begin{center}\large#1\end{center}
    \vskip0.5em}%
}

\title{$\mathbb{MENRVA}$}
\subtitle{@futurice freaklies}
\author{Oleg Grenrus}
\date{July 03, 2014}

\DeclareMathOperator*{\inr}{\mathsf{inr}}
\DeclareMathOperator*{\inl}{\mathsf{inl}}
\DeclareMathOperator*{\none}{\mathsf{none}}
\DeclareMathOperator*{\some}{\mathsf{some}}
\DeclareMathOperator*{\ap}{\mathsf{ap}}
\DeclareMathOperator*{\fst}{\mathsf{fst}}
\DeclareMathOperator*{\snd}{\mathsf{snd}}

% Theorems
% http://en.wikibooks.org/wiki/LaTeX/Theorems

\theoremstyle{definition}
\newtheorem{definition}{Definition}[section]
\newtheorem{example}[definition]{Example}
\newtheorem{exercise}[definition]{Exercise}
\newtheorem{note}[definition]{Note}

\theoremstyle{plain}
\newtheorem{proposition}[definition]{Proposition}
\newtheorem{lemma}[definition]{Lemma}
\newtheorem{theorem}[definition]{Theorem}

\pagestyle{empty}

\definecolor{blue}{rgb}{0,0,0.5}
\definecolor{green}{rgb}{0,0.5,0}
\definecolor{red}{rgb}{0.5,0,0}
\definecolor{orange}{rgb}{1,0.5,0}
\definecolor{darkorange}{rgb}{0.5,0.25,0}
\definecolor{lightgrey}{rgb}{0.7,0.7,0.7}
\definecolor{grey}{rgb}{0.5,0.5,0.5}

\setlength\parindent{0pt}
\setlength{\parskip}{3mm plus1mm minus1mm}

\newcommand{\identifier}[1]{\texttt{\color{blue}#1}}
\newcommand{\constructor}[1]{\ensuremath{\mathsf{\color{green}#1}}}
\newcommand{\functor}[1]{\ensuremath{\mathcal{\color{darkorange}#1}}}

\begin{document}


\maketitle
\thispagestyle{empty}
\newpage

{\LARGE$\mathbb{MENRVA}$-- yet another reactive library}

\href{https://github.com/baconjs/bacon.js}{\textsc{Bacon.js}} is (almost) non-glitchy, but not-so performant.
\href{https://github.com/Reactive-Extensions/RxJS}{\textsc{RxJS}} is glitchy but performant.

I begin to bias towards using cold-observables of \textsc{RxJS}\footnote{Can use \textsc{Bacon.js} too, but I don't need glitch-free events} (the good parts?) with data-flow propagation of something else ($=$ \textsc{Menrva}).

\begin{center}
\includegraphics[scale=0.7]{figures/data-events.mps}
\end{center}

\newpage

\section*{Semantics}

Following definitions are from \emph{Push-pull functional reactive programming}\footnote{Conal Elliott, \emph{Push-pull functional reactive programming}, Haskell Symposium, 2009, \url{http://conal.net/papers/push-pull-frp/}}.

\subsection*{FRP -- Behavior}

Semantically, a \emph{(reactive) behavior} is just a function of time. In \textsc{Menrva} we call them \emph{signals}.
\begin{align*}
\mathcal{S}\,\alpha = \mathcal{T} \to \alpha
\end{align*}

Historically in FRP, $\mathcal{T} = \mathbb{R}$. As we’ll see, however, the semantics of behaviors assumes only that $\mathcal{T}$ is totally ordered. 

\newpage

\subsubsection*{Applicative functor}

Pure functions can be "lifted" to apply to signals. Classic FRP (CFRP) had a family of lifting combinators:
\begin{align*}
\texttt{lift}_n &:: (\alpha_1 \to \cdots \alpha_n \to \beta) \to (\mathcal{S}\,\alpha_1 \to \cdots \mathcal{S}\,\alpha_n \to \mathcal{S}\,\beta)
\end{align*}
Lifting is pointwise and synchronous:
\begin{align*}
\texttt{lift}_n\,f\,b_1\cdots b_n &= \lambda\,t \mapsto f (b_1\, t)\cdots (b_n\, t)
\end{align*}
The \emph{Functor} \texttt{map} is $\texttt{lift}_1$:
\begin{align*}
\texttt{map}\,f\,b = f \circ b
\end{align*}

\newpage

In \textsc{Menrva} \texttt{map} is \texttt{map} and $\texttt{lift}_n$ is \texttt{combine} (with function parameter at the end):
\begin{verbatim}
var $sq = $source.map(function (x) {
  return x * x;
});

var $quad = menrva.combine($a, $b, $c, function (a, b, c) {
  return a * b + c;
});
\end{verbatim}

And that's the whole (functional) API. 

{\small There is \texttt{onValue} to observe changes. And \texttt{set} and \texttt{modify} to initiate changes.}

\newpage

\subsubsection*{Monad}

Although Signal is a semantic \emph{Monad} as well, \textsc{Menrva} doesn't have monadic \texttt{join} combinator.

Semantic of \texttt{join} is simple, but the correct implementation hard.
\begin{align*}
\texttt{join}\,s = t \mapsto s\,t\,t
\end{align*}

Monad instance would make possible to make dynamic data flow networks. With applicative functor the shape of the network is static (and you can't make loops) $\equiv$ simple and easy.

\newpage

\section*{Examples}

\url{https://github.com/phadej/menrva/tree/master/examples}

\begin{itemize}
\item \emph{counter} examples (pun not-intended). With data-first approach, you segregate \emph{query} from \emph{control}.
\item \emph{suggestions} \href{https://gist.github.com/staltz/868e7e9bc2a7b8c1f754}{\emph{The introduction to Reactive Programming you've been missing}}
\end{itemize}

\newpage


\section*{What are plans for \textsc{Menrva}?}

Write more examples, to verify that \textsc{Menrva}'s simplistic approach is enough for the real world.

\hspace{5mm}\vdots

Use in the real world application.

\hspace{5mm}\vdots

PROFIT!

\newpage

\section*{What are you going to add still?}

Something to handle network destruction.
% Currently all nodes are double linked (parent $\leftrightarrow$ child).
Thus dynamically created (and destroyed) UI components create leaf-subnetwork which should be cleaned up.
\textsc{JavaScript} doesn't have weak references (which might be enough to solve the problem).
So we need some manual mechanism for cleanup.

One idea is to make own reference counting:
\begin{verbatim}
var $state = menrva.source(...);
var $derived = $state.map(...).retain();

var $componentData = $derived.map(...).retain();
// cleanup:
$componentData.release(); // also removes all callbacks etc.
\end{verbatim}
\newpage


\section*{What are you going to add still? -- 2}

Should we have monadic join? (= \texttt{flatMap}). Do we really need it? Can be worked around, but sometimes it is convenient:
{\small
\begin{verbatim}
$artistImage = $artist.flatmap(getArtistImageSignal);

function getArtistImageSignal(artist) {
  if (!cache[artist.id]) { 
    cache[artist.id] = menrva.source(placeholderImage);
    
    // schedule carousel for images...
  }
  return cache[artist.id];
}
\end{verbatim}
}
Here the cache (or artist image store) can be signal itself -- no need for \texttt{flatMap}.


\section*{What are you going to add still? -- 3}

\emph{Borrow} goodies from \textsc{Bacon.js} -- Function construction rules:
\begin{verbatim}
$signal.map(".foo"); // shorter way of saying:
$signal.map(function (x) {
  return x.foo;
});
\end{verbatim}

\newpage


\section*{What are you going to add still? -- 4}

Lenses. Sub-value of source acting as source itself.
\begin{verbatim}
var $foobar = $source.lens("foo.bar");

$foobar.set(tx, value); // the same as
$source.modify(tx, function (source) {
  return assocIn(source, ["foo", "bar"], value);
});
\end{verbatim}
\texttt{assocIn} -- \url{http://clojuredocs.org/clojure_core/clojure.core/assoc-in}

\newpage

\section*{And why the choice for transactions?}

You able to update many source simultaneously, avoiding immediate propagation:
\begin{verbatim}
var $a = menrva.source(1);
var $b = menrva.source(2);

menrva.combine($a, $b, add).onValue(console.log);

var tx = menrva.transaction();
$a.set(tx, 2);
$b.set(tx, 3);
tx.commit();
\end{verbatim}

\section*{And why the choice for transactions? -- 2}

Also to mention, transactions arise naturally, if you try to implement anything similar in
\textsc{Haskell} or \textsc{Clojure} (using \textsc{STM}).

There you will be forced to do all mutation inside STM transaction.

Changes initiated in transaction are applied only when you commit the transaction: you \emph{always} have consistent data.

\newpage
{\LARGE the end}
\newpage


\end{document}
